\subhead{摘要}
即使顶棚排烟口能够捕集火源上方的烟气，狭长房间远端的烟层仍可能下降。本研究使用火灾动力学模拟软件（FDS），考察侧开口对顶棚射流温度的影响，并确定维持烟层净高所需的排烟量。与丙烷火实验对照发现，FDS 捕捉到了温度峰值向封闭侧偏移的趋势，温度误差的大小则随门态和火源功率而变化。针对丙烯火，我们在 16 个已有的挡烟垂壁与排烟量组合上，分别回放贝叶斯优化（BO）和遗传算法（GA）的搜索过程。两种算法均比固定顺序更早找到已测方案中净高余量最大的组合。随后，我们评估了改进布局，在顶棚设置两个排烟口，并按门态设定端墙补风口的高度。逐一用 FDS 计算六种门态与烟炱产率组合耗时较长，因此我们引入物理信息神经网络（PINN）快速预测各工况的最低烟层高度，并用已有 FDS 结果回放这一预筛策略，检验优先检查预测净高较低工况的效果。细网格 FDS 结果表明，1.12~m$^3$/s 是所测风量中满足约束且估算功率最低的选择。在六个工况各自的 120~s 模拟中，八个测点的最低烟层高度比 0.38~m 目标高出 55.8~mm。

\reportheading{1.\quad 引言}
火羽流撞击顶棚后，热气体沿顶棚向两侧流动。自然通风隧道的缩尺实验沿这一流动方向测量温度，记录顶棚射流随距离的衰减~\cite{tunnel}。在一侧设有多个门洞的狭长房间内，还需考虑补风从哪一侧进入。Peng 等~\cite{peng}在这一构型下改变侧门的开启组合，测量了顶棚温升沿房间长度的分布。他们的实验表明，羽流附近的温度取决于具体开启哪些门，随开口数量增加并无简单的变化规律。因此，沿房间纵向测得的温度剖面，可以同时检验模型对峰值温升及其位置的预测。

任务 1 利用 Peng 等的丙烷火实验，在四门全开和仅 R1 开启两种条件下检验上述特征。门口速度与温度剖面可用来判断冷空气从何处进入火源区、热气体从何处排出。顶棚机械排烟时，补充空气也沿这些通路进入房间。因此，任务 2 在两种门态下设计丙烯火排烟系统，针对不同门态预设端墙补风口，使两种门态与三种烟炱产率组合成的六个工况均保有 0.38~m 净高，并在所测风量中寻找功率最低者。我们回放原布局的设计搜索，并在改进布局中考察风量与工况的检查次序如何影响查询次数。风量推荐以 FDS 记录的最低烟层高度为依据。
